\scanpage{84}
% archive sheet number: 83
% block display | line 1
\[
P\bigl(\sum f(\unsure{\gamma|\sigma})\cdot\gamma a\bigr)\ne0
\]
% block text | lines 2-9
où $P$ est un polynôme à coeff. dans $K$ \deleted{défini} défini sur l'espace $K$-vectoriel $A$, \unsure{p.ex. en des} coordonnées \deleted{\unsure{$a$}} du point courant, p.ex. à l'aide d'une base. On montre d'abord qu'on peut trouver une base formée d'éléments invariants par $\Gamma$ [c'est un lemme, qui n'utilise que la structure \deleted{\unsure{d'algèbre}}\inserted{$K$-$\Gamma$-espace} vectoriel de $A$, -- on n'utilise pas que la dimension soit finie]

% block text | line 10
Si $a=(a_i)_{1\leq i\leq\unsure{n}}$, on a donc
% block display | line 11 cancelled
\[
\deleted{P\bigl(\sum f(\gamma)\gamma a\bigr)=\unsure{Q|R}\bigl(f_1(\gamma),\ldots,f_n(\gamma)\bigr);}
\]
% block margin note | right of lines 10-11
% slanted
\margin{$R$ un polynôme en $n|\Gamma|$ variables}
% block display | line 12
\[
P\bigl(\sum f(\gamma)\gamma(a_\gamma)\bigr)=R\bigl(\{\{\gamma a_\gamma^1,\ldots,\gamma a_\gamma^n\}\}_{\gamma\in\Gamma}\bigr)
\]
% block text | lines 13-17
Supposons $K$ infini. Si l'\unsure{équation|expression} précédente était nulle \inserted{en y faisant $a_\gamma^i=a^i$ pour $\forall\gamma$ (et tout $i$)} quels que soient $a^1,\ldots,a^n\in K$, alors d'après le th. d'indépendance algébrique des automorphismes, alors $R$ est identiquement nul, en particulier on a \deleted{\unsure{$\sum$}} aussi
% block display | line 18
\[
\deleted{\unsure{R(\{\gamma a_\gamma^i}}\qquad P\bigl(\sum f(\gamma)\gamma(a_\gamma)\bigr)=0
\]
% block text | lines 19-20
quel que soit le choix des $a_\gamma$, d'où (faisant $a_{\gamma_0}=1$ et $\unsure{a_{\gamma_0}|a_\gamma}=0$ si $\gamma\ne\gamma_0$)
% block display | line 21
\[
P(f(\gamma_0))=0,\quad\text{absurde.}
\]
% block text | lines 21-23
/ Dans le cas où $K$ est fini, il faut utiliser \unsure{les résultats|le résultat} de Lang.....
